% !TeX program = xelatex
% !TeX encoding = UTF-8
% 实验报告模板：个人配置 -> 排版设置 -> 常用组件 -> 正文示例。
% 从三个报告提炼：A4 / 12pt / 四边 2.5cm / 1.35 倍行距。
% 每次使用请复制整个文件夹；用 XeLaTeX 编译两次以更新目录和引用。
\documentclass[UTF8,a4paper,12pt,fontset=fandol]{ctexart}

\usepackage{geometry}
\usepackage{graphicx}
\usepackage{float}
\usepackage{listings}
\usepackage{xcolor}
\usepackage{setspace}
\usepackage{enumitem}
\usepackage{amsmath,amssymb}
\usepackage{booktabs,tabularx}
\usepackage[hidelinks]{hyperref}

% ======================================================================
% 1. 个人配置：修改已有定义的右侧内容，不要重复定义同名命令。
% ======================================================================
\newcommand{\ReportTitle}{某某课程实验报告}
\newcommand{\TopicLabel}{实验题目} % 可改为“实验选题”或“课程设计题目”
\newcommand{\ReportTopic}{在此填写实验或课程设计题目}
\newcommand{\StudentName}{你的姓名}
\newcommand{\StudentID}{你的学号}
\newcommand{\Supervisor}{指导教师姓名}
\newcommand{\ReportDate}{2026年10月5日} % 按实际完成日期修改
\newcommand{\Institution}{xx大学计算机科学与技术学院}
\newcommand{\LogoFile}{images/logo.png} % 可选：放入自己的校徽；缺图时显示占位符
% 不需要仓库地址时留空；URL 中的下划线不需要写成 \_。
\newcommand{\RepositoryURL}{}
\newcommand{\SummaryLabel}{摘要} % 也可改为“自我评价”
\newcommand{\ReportSummary}{在此简述实验对象、实现方法、验证方式与主要结论。建议控制在一百至一百五十字；不需要时关闭封面摘要。}

\newif\ifShowLogo
\ShowLogotrue
\newif\ifShowTopic
\ShowTopictrue
\newif\ifShowSummary
\ShowSummaryfalse % 改为 \ShowSummarytrue 显示封面摘要
\newif\ifShowContents
\ShowContentstrue

% ======================================================================
% 2. 排版设置：通常无需修改。
% 显式使用 Fandol 中文字体；TeX Live 自带，无需上传 fonts/。
% ======================================================================
\geometry{left=2.5cm,right=2.5cm,top=2.5cm,bottom=2.5cm}
\setstretch{1.35}
\setlength{\parindent}{2em}
\setcounter{tocdepth}{2} % 目录显示到 subsection
\setcounter{secnumdepth}{3}
\urlstyle{same}
\hypersetup{pdftitle={\ReportTitle},pdfauthor={\StudentName}}
\setlist[itemize]{leftmargin=2em,itemsep=0.2em,topsep=0.4em}
\setlist[enumerate]{leftmargin=2em,itemsep=0.2em,topsep=0.4em}

\lstset{
  basicstyle=\ttfamily\small,
  breaklines=true,
  frame=single,
  numbers=left,
  numberstyle=\tiny,
  numbersep=8pt,
  xleftmargin=1.8em,
  framexleftmargin=0.5em,
  tabsize=4,
  showstringspaces=false,
  columns=fullflexible,
  keepspaces=true,
  aboveskip=0.8em,
  belowskip=0.6em
}

% ======================================================================
% 3. 常用组件。
% ReportFigure[最大高度]{图片路径}{图注}{引用标签}。
% 缺图时显示占位框：模板本身不依赖任何实验截图。
% 图片宽度上限为当前行宽的 90%，保持原比例，不拉伸。
% ======================================================================
\newcommand{\ReportImage}[2][6cm]{%
  \IfFileExists{#2}{%
    \includegraphics[width=0.9\linewidth,height=#1,keepaspectratio]{#2}%
  }{%
    \fbox{%
      \begin{minipage}[c][#1][c]{0.85\linewidth}
        \centering\small
        图片占位：请将实验截图放入 images/。\par
        \medskip\texttt{\detokenize{#2}}
      \end{minipage}%
    }%
  }%
}
\newcommand{\ReportFigure}[4][6cm]{%
  \begin{figure}[H]
    \centering
    \ReportImage[#1]{#2}
    \caption{#3}
    \label{#4}
  \end{figure}%
}
\newcommand{\ReportNote}[2]{%
  \par\medskip\noindent
  \fbox{\begin{minipage}{\dimexpr\linewidth-2\fboxsep-2\fboxrule\relax}
    \textbf{#1}\par #2
  \end{minipage}}\par\medskip
}

% 封面中的下划线信息栏，仅适合较短的姓名、学号、日期。
\newcommand{\CoverField}[2]{%
  #1：\underline{\makebox[8cm][c]{#2}}\par\vspace{0.35cm}%
}
\newcommand{\MakeReportCover}{%
  \begin{titlepage}
    \centering
    \vspace*{0.8cm}
    \ifShowLogo
      \IfFileExists{\LogoFile}{%
        \includegraphics[width=0.8\linewidth,height=2.2cm,keepaspectratio]{\LogoFile}\par
      }{%
        \fbox{\parbox[c][2.2cm][c]{0.75\linewidth}{%
          \centering\zihao{4}\heiti 校徽占位符\par
          \medskip\small\ttfamily images/logo.png
        }}\par
      }
      \vspace{1.2cm}
    \fi
    {\zihao{2}\heiti\ReportTitle\par}
    \vspace{1.5cm}
    \ifShowTopic
      {\zihao{3}\heiti\TopicLabel：\par}
      \vspace{0.4cm}
      {\zihao{3}\heiti\ReportTopic\par}
      \vspace{1.4cm}
    \fi
    {\zihao{4}
      \CoverField{姓名}{\StudentName}
      \CoverField{学号}{\StudentID}
      \CoverField{指导教师}{\Supervisor}
      \CoverField{完成日期}{\ReportDate}
    }
    % 不用固定宽度下划线包住 URL，长地址可自行换行。
    \ifx\RepositoryURL\empty\else
      \begin{minipage}{0.95\linewidth}
        \centering\small 仓库地址：\expandafter\url\expandafter{\RepositoryURL}
      \end{minipage}\par
    \fi
    \vfill
    \ifShowSummary
      \begin{minipage}{\linewidth}
        \setlength{\parindent}{2em}
        \textbf{\SummaryLabel：}\par\ReportSummary
      \end{minipage}\par\vspace{0.8cm}
    \fi
    {\zihao{4}\Institution\par}
    \vspace*{0.4cm}
  \end{titlepage}%
}

% ======================================================================
% 4. 正文：从这里开始替换；示例文字、结果占位和文献必须按实际填写。
% ======================================================================
\begin{document}
\MakeReportCover

% 封面无页码；目录用罗马数字；正文从阿拉伯数字 1 开始。
\ifShowContents
  \pagenumbering{Roman}
  \tableofcontents
  \clearpage
\fi
\pagenumbering{arabic}

\section{实验目的与任务}
说明实验背景、待解决的问题以及本次需要完成的任务。课程设计可将本节改为“研究背景与意义”，再增加“研究内容”。
\begin{enumerate}
  \item 掌握本实验涉及的关键原理与方法。
  \item 完成核心模块的设计、实现与运行。
  \item 验证结果正确性，并分析实验现象或性能。
\end{enumerate}

\section{实验环境}
记录能复现实验的硬件、软件与版本信息；区分本机环境和服务器环境。
\begin{table}[H]
  \centering
  \caption{实验环境配置（待填写）}
  \label{tab:environment}
  \renewcommand{\arraystretch}{1.3}
  \begin{tabularx}{0.95\linewidth}{l X}
    \toprule
    项目 & 配置说明 \\
    \midrule
    硬件平台 & 填写处理器、内存或开发板型号 \\
    操作系统 & 填写系统名称与版本 \\
    开发工具 & 填写编辑器、编译器或仿真器及版本 \\
    编程语言 & 填写语言名称与标准版本 \\
    构建与运行 & 填写构建方式、关键参数或依赖 \\
    \bottomrule
  \end{tabularx}
\end{table}

\section{总体设计与关键原理}
\subsection{总体方案}
说明输入、输出、模块划分与模块之间的关系。项目结构可用无语言高亮的代码块展示。
\begin{lstlisting}[numbers=none]
project/
|-- src/
|-- include/
|-- tests/
`-- results/
\end{lstlisting}

\subsection{关键原理}
解释算法、体系结构或硬件机制，并给出必要的公式。正文引用的资料见文献~\cite{ref:material}。
\begin{equation}
  S = \frac{T_{\mathrm{baseline}}}{T_{\mathrm{optimized}}}
  \label{eq:speedup}
\end{equation}
其中，$T_{\mathrm{baseline}}$ 和 $T_{\mathrm{optimized}}$ 分别表示基线版本与优化版本的运行时间。若实验不涉及性能比较，请换成适用的原理公式。

\section{实验过程与实现}
\subsection{环境准备与运行方法}
说明依赖、构建步骤与运行命令。以下命令仅演示排版，请换成项目实际命令。
\begin{lstlisting}[numbers=none]
make
./program --input test.dat
\end{lstlisting}

\subsection{实验一：核心功能实现}
\subsubsection{实验要求}
写明需要实现的功能、输入条件与验收标准。

\subsubsection{实现过程}
说明关键修改和原因；只展示与论述有关的代码。中文说明建议放在代码块外。
\begin{lstlisting}[language=C++,caption={代码排版示例},label={lst:example}]
int sum(const int* data, int n) {
    int result = 0;
    for (int i = 0; i < n; ++i) {
        result += data[i];
    }
    return result;
}
\end{lstlisting}

\subsubsection{验证与结果}
记录测试输入、预期输出、实际输出与判断依据。图~\ref{fig:verification} 应替换为本次实验的截图、波形或结果图。
\ReportFigure[5cm]{images/verification.png}{实验验证结果（图片待替换）}{fig:verification}

% 多模块实训：复制上述 subsection 及三个 subsubsection，依次填实验二、三……。
% 若每个实验内容较多，可将“实验一”提升为 section，其内部标题改为 subsection。

\section{实验结果与分析}
\subsection{结果记录}
先给出实际记录，再说明结论。以下表格只提供结构，不包含任何实验数据。
\begin{table}[H]
  \centering
  \caption{实验结果记录（待填写）}
  \label{tab:results}
  \renewcommand{\arraystretch}{1.3}
  \begin{tabularx}{0.95\linewidth}{l X X X}
    \toprule
    测试条件 & 预期结果 & 实际结果 & 验证结论 \\
    \midrule
    条件一 & 待填写 & 待填写 & 待填写 \\
    条件二 & 待填写 & 待填写 & 待填写 \\
    \bottomrule
  \end{tabularx}
\end{table}

\subsection{现象与原因分析}
结合表~\ref{tab:results} 分析实际结果是否符合预期，并说明证据与限制。性能实验可使用公式~\eqref{eq:speedup}，同时明确计时范围、测试次数与统计方式；硬件实验可分析波形与时序。

\ReportNote{写作提示}{截图下方应解释“观察到了什么、说明了什么”。没有验证的数据和推测应明确标注。正式报告可删除本提示框。}

\section{问题与解决方法}
\subsection{问题一：在此填写问题名称}
\begin{description}[leftmargin=2em,style=nextline]
  \item[现象] 描述可观察到的错误及触发条件。
  \item[原因] 说明定位过程及根本原因。
  \item[解决方法] 记录修改内容与修复后的验证结果。
\end{description}

\section{总结与展望}
总结实际完成的工作、已验证的结果和掌握的方法；说明尚未解决的问题及后续改进方向。避免直接保留模板中的示例描述。

% 参考文献与 Beamer 模板用法相同：bibitem 定义键，cite 使用键。
% 无需 .bib 文件、BibTeX 或 Biber。替换示例条目，保留唯一键名。
\clearpage
\phantomsection
\addcontentsline{toc}{section}{参考文献}
\begin{thebibliography}{99}
  \bibitem{ref:material}
    作者或机构。\emph{资料名称（待替换）}。版本或出版信息，年份。
\end{thebibliography}

% 可选附录：取消注释后填写，适合补充测试记录或较长代码。
% \clearpage
% \appendix
% \section{补充实验记录}
% 在此填写正文之外的补充材料。
\end{document}
